\begin{frame}{``보상을 잘 설계하는 것'' 자체가 공학}
  지도학습의 라벨은 \textbf{외부에서 주어짐} --- RL의 보상은
  \textbf{내가 설계}. 보상을 잘못 설계하면 에이전트는
  \kq{보상이 원하는 것이 아니라 \textbf{보상을 회피하는 방법}을 배운다}.
  \vskip0.6em
  \kb{보상 해킹(reward hacking)}의 고전 두 사례:
  \begin{itemize}
    \item 진흙탕 로봇, ``진흙 밟기 $-1$'' $\to$ 진흙을 \textbf{지나는 대신
      보지 못하게}(센서 덮기) 하는 전략.
    \item Atari \textit{Enduro}: ``점수를 안 잃는 것''을 보상 $\to$
      \textbf{차를 도로 가장자리에 붙여 멈추기}. (``이기는 법''은 아니지만
      보상은 더 이상 안 잃음.)
  \end{itemize}
  \vskip0.4em
  \textbf{보상 설계 = 목적을 정밀하게 언어화하는 작업}.
  13$\sim$14주차에서 재등장 --- 실제 로봇 보상(속도$\cdot$토크$\cdot$자세)의
  각 계수가 성과의 80\%를 결정.
\end{frame}
